%
% 第4章：神经网络架构

\clearpage
\cleardoublepage

\chapter{神经网络架构}

\section{引言}

神经网络架构定义了神经元的组织方式和连接方式，以处理信息。本章介绍三种基本神经网络架构：深度神经网络(DNN)、卷积神经网络(CNN)和循环神经网络(RNN)。

\section{深度神经网络（DNN）}

深度神经网络由多层神经元组成，信息从输入流向输出，无环路。

全连接DNN执行：
\begin{equation}
\mathbf{h}^{(l)} = f^{(l)}(\mathbf{W}^{(l)}\mathbf{h}^{(l-1)} + \mathbf{b}^{(l)})
\end{equation}

属性：
\begin{itemize}
    \item \textbf{通用逼近：} 可以逼近任何连续函数
    \item \textbf{适合：} 表格数据、基于特征的分类
\end{itemize}

\section{卷积神经网络（CNN）}

CNN通过局部连接和参数共享利用数据的空间结构。

卷积操作：
\begin{equation}
y[i,j] = \sum_{m}\sum_{n} w[m,n] \cdot x[i+m, j+n]
\end{equation}

关键参数：
\begin{itemize}
    \item \textbf{卷积核大小：} $K \times K$（通常 $3 \times 3$, $5 \times 5$）
    \item \textbf{步长：} $s$（步长，通常1或2）
    \item \textbf{填充：} $p$（零填充控制输出大小）
\end{itemize}

输出大小：
\begin{equation}
H_{out} = \left\lfloor \frac{H_{in} + 2p - K}{s} + 1 \right\rfloor
\end{equation}

属性：
\begin{itemize}
    \item \textbf{局部感受野}
    \item \textbf{参数共享}
    \item \textbf{平移等变性}
    \item \textbf{适合：} 图像、视频、空间数据
\end{itemize}

\section{循环神经网络（RNN）}

RNN通过维护隐藏状态处理序列数据。

对于输入序列 $\mathbf{x}^{(1)}, \mathbf{x}^{(2)}, \ldots, \mathbf{x}^{(T)}$：
\begin{align}
\mathbf{h}^{(t)} &= f(\mathbf{W}_{xh}\mathbf{x}^{(t)} + \mathbf{W}_{hh}\mathbf{h}^{(t-1)} + \mathbf{b}_h) \\
\mathbf{y}^{(t)} &= g(\mathbf{W}_{hy}\mathbf{h}^{(t)} + \mathbf{b}_y)
\end{align}

属性：
\begin{itemize}
    \item \textbf{序列处理}
    \item \textbf{共享参数}
    \item \textbf{记忆：} 隐藏状态捕获过去信息
    \item \textbf{适合：} 文本、时间序列、音频
\end{itemize}

\section{长短期记忆网络（LSTM）}

LSTM通过门控机制解决梯度消失问题。

LSTM门：
\begin{align}
\mathbf{f}^{(t)} &= \sigma(\mathbf{W}_f[\mathbf{h}^{(t-1)}, \mathbf{x}^{(t)}]) \quad \text{(遗忘门)} \\
\mathbf{i}^{(t)} &= \sigma(\mathbf{W}_i[\mathbf{h}^{(t-1)}, \mathbf{x}^{(t)}]) \quad \text{(输入门)} \\
\mathbf{o}^{(t)} &= \sigma(\mathbf{W}_o[\mathbf{h}^{(t-1)}, \mathbf{x}^{(t)}]) \quad \text{(输出门)}
\end{align}

属性：
\begin{itemize}
    \item \textbf{选择性记忆}
    \item \textbf{长期依赖}
    \item \textbf{广泛使用}
\end{itemize}

\section{门控循环单元（GRU）}

GRU用更少的门简化了LSTM：
\begin{align}
\mathbf{z}^{(t)} &= \sigma(\mathbf{W}_z[\mathbf{h}^{(t-1)}, \mathbf{x}^{(t)}]) \quad \text{(更新门)} \\
\mathbf{r}^{(t)} &= \sigma(\mathbf{W}_r[\mathbf{h}^{(t-1)}, \mathbf{x}^{(t)}]) \quad \text{(重置门)}
\end{align}

比LSTM参数更少，性能通常相当。

\section{架构选择}

\begin{table}[h]
    \centering
    \caption{架构选择指南}
    \begin{tabular}{L{2.5cm} L{3.5cm} L{4cm}}
        \toprule
        \textbf{架构} & \textbf{最佳用途} & \textbf{示例} \\
        \midrule
        DNN & 表格数据 & 分类、回归 \\
        CNN & 图像、空间数据 & 目标检测、分割 \\
        RNN/LSTM & 序列 & 自然语言处理、时间序列 \\
        Transformer & 长序列 & 翻译、生成 \\
        \bottomrule
    \end{tabular}
\end{table}

\section{本章小结}

\begin{enumerate}
    \item DNN是用于通用学习的全连接网络
    \item CNN通过局部连接利用空间结构
    \item RNN通过循环连接处理序列数据
    \item LSTM和GRU通过门控解决梯度消失问题
    \item 架构选择取决于数据类型和任务
\end{enumerate}
